\chapter{Thiết kế và cài đặt hệ thống}\label{chap:implementation}

\section{Kiến trúc xử lý}
Lớp COKB nhận dữ liệu JSON chuẩn hóa hoặc dữ liệu chuyển đổi từ Oxigraph.
Với lệnh build/reason/import, CLI kiểm tra cấu trúc trước, chỉ suy luận
khi không có lỗi mức error, sau đó xuất store cùng các graph và proof.
SHACL được gọi bằng lệnh riêng sau khi xuất RDF.

\begin{figure}[htbp]
\centering
\begin{tikzpicture}[node distance=0.65cm and 1.2cm,
  block/.append style={text width=4.5cm}]
  \node[block] (input) {JSON mẫu\\Adapter Oxigraph};
  \node[block,below=of input] (validation) {Kiểm tra Python\\giá trị bắt buộc, trùng ID};
  \node[block,right=of validation,text width=3.4cm] (quarantine) {Có lỗi: báo cáo\\và quarantine};
  \node[block,below=of validation] (engine) {Operators + Rules\\đánh giá có thứ tự};
  \node[block,below=of engine] (export) {Store JSON và ba file RDF\\Asserted / Inferred / Proof};
  \node[block,below=of export] (query) {Query / Explain\\Assessment và bằng chứng};
  \node[block,right=of export,text width=3.4cm] (shacl) {SHACL tùy chọn\\kiểm tra RDF đã xuất};
  \draw[flow] (input)--(validation);
  \draw[flow] (validation)--node[above,font=\scriptsize]{error}(quarantine);
  \draw[flow] (validation)--node[left,font=\scriptsize]{hợp lệ}(engine);
  \draw[flow] (engine)--(export);
  \draw[flow] (export)--(query);
  \draw[flow] (export)--(shacl);
\end{tikzpicture}
\caption{Luồng COKB được thực thi bởi các lệnh CLI hiện tại.}
\label{fig:architecture}
\end{figure}

\section{Phân rã module}
\begin{table}[htbp]
\centering\small
\caption{Trách nhiệm các module của lớp COKB}\label{tab:modules}
\begin{tabularx}{\textwidth}{p{4.0cm}Y}
\toprule
Module trong \src{scripts/cokb/} & Trách nhiệm\\
\midrule
\src{model.py} & Dataclass, enum mức đánh giá, ID ổn định, validation report.\\
\src{operators.py} & Chuẩn hóa, tập từ, xung đột, điều kiện so sánh.\\
\src{rules.py}, \src{reasoner.py} & Điều kiện luật, chọn luật, dựng assessment và proof.\\
\src{validation.py}, \src{shacl.py} & Kiểm tra Python và kiểm tra RDF bằng pySHACL.\\
\src{repository.py} & Nạp/lưu, truy vấn, giải thích, tuần tự hóa RDF và JSON.\\
\src{importer.py} & Đọc Oxigraph ở chế độ read-only và chuyển cặp mã.\\
\src{xlsx.py}, \src{evaluation.py} & Đọc sheet gold bằng thư viện chuẩn, tính chỉ số đánh giá.\\
\src{cli.py} & Khai báo subcommand, kiểm soát luồng, mã thoát, xuất báo cáo.\\
\bottomrule
\end{tabularx}
\end{table}

Lõi suy luận và đọc XLSX có thể chạy bằng thư viện chuẩn Python.
Theo \src{pyproject.toml}, môi trường project yêu cầu Python từ 3.11 đến
trước 3.13; các phụ thuộc tùy chọn nhóm validation là
\src{rdflib} và \src{pyshacl}. Lệnh import cần \src{pyoxigraph}.
Các chức năng COKB được gọi từ \src{main.py} mà không phải cấu hình LLM.

\section{Định danh và nhập dữ liệu}
Nếu input chưa cung cấp ID, \src{stable_mapping_id} băm chuỗi nối hệ mã,
mã nguồn, hệ mã và mã đích bằng SHA-256, lấy 16 ký tự hex đầu:
\begin{equation}
 id(m)=\texttt{mapping-}\ \Vert\
 \operatorname{prefix}_{16}\!\left(
 \operatorname{SHA256}(s_s\Vert | \Vert k_s\Vert | \Vert s_t\Vert | \Vert k_t)
 \right).
\end{equation}
Định danh này bỏ qua label và evidence. Hai mapping có cùng bốn thành
phần nhưng nguồn khác nhau có thể cùng ID; đây là giới hạn thiết kế,
không phải cơ chế quản lý lịch sử nhiều bằng chứng.

Adapter dùng SPARQL \texttt{SELECT DISTINCT} trong các named graph để lấy
cặp mã và nhãn, rồi khử trùng theo bộ $(s_s,k_s,s_t,k_t)$.
Các tham số \texttt{source-system}, \texttt{target-system} và
\texttt{evidence-source} được truyền từ CLI, không được suy ra riêng cho
từng dòng. Vì vậy, nếu nguồn chứa lẫn ICD10WHO và ICD10CM thì việc gán
một hệ nguồn cho toàn bộ kết quả có thể làm sai metadata.
Code cũng chưa bảo toàn graph IRI và mapping IRI của từng dòng làm provenance.

\section{Xử lý lỗi và các artifact}
Khi validation phát hiện error, CLI ghi
\src{validation_report.json} và \src{quarantine.json}, trả mã thoát 1 và
không tiếp tục inference của lần gọi đó. Cả batch bị dừng; không có cơ chế
tự động xử lý phần hợp lệ còn lại. API \src{infer} và \src{explain}
không tự bắt buộc validation, nên ứng dụng tích hợp phải gọi kiểm tra rõ ràng.

Khi thành công, \src{save} tạo sáu artifact:
\src{store.json}, \src{validation_report.json}, \src{proofs.json},
\src{asserted_graph.ttl}, \src{inferred_graph.ttl}, \src{proof_graph.ttl}.
SHACL hợp nhất ba file RDF cùng ontology để kiểm tra; tùy chọn suy luận
RDFS trong bước này phục vụ nhận biết lớp và ràng buộc,
không thay thế bộ đánh giá A/B/C bằng Python.

\placeholderfigure{images/cli_demo.png}
{Kết quả truy vấn COKB cho mã nguồn K05.1.}
{Chụp terminal sau khi chạy cokb\_build và cokb\_query.
Hiển thị hai target DA0B.Y và DA0Z, mức A và UNASSESSED,
cùng mã luật của mỗi assessment. Chỉ dùng output thật của repo.}
{fig:cli}

\section{Ví dụ giải thích một mapping}
Với cặp K05.1 và DA0B.Y trong fixture, hai nhãn đều là
\emph{Chronic gingivitis}. Sau chuẩn hóa, cả hai thành
\texttt{chronic gingivitis}. Luật đầu tiên khớp, trả A với hai premise
là nhãn đã chuẩn hóa và nguồn do fixture khai báo.
\begin{lstlisting}[caption={Trích phần proof của ví dụ nhãn trùng nhau},label={lst:proof}]
{
  "conclusion": {
    "mapping_id": "mapping-k05-1-da0b-y",
    "predicate": "hasMappingLevel",
    "object": "A"
  },
  "rule_id": "R-EXACT-LABEL",
  "sources": ["icd10cm_to_icd11_2024_full"]
}
\end{lstlisting}
Mã minh họa~\ref{lst:proof} lược bỏ mảng premises cho gọn.
Chuỗi source trong fixture là metadata tự khai báo; tên của nó không
chứng minh ví dụ đã được xác minh với toàn bộ file ICD.

\placeholderfigure{images/proof_graph.png}
{Trực quan hóa vết suy luận cho mapping K05.1 đến DA0B.Y.}
{Nạp ba file Turtle và ontology vào công cụ xem RDF.
Vẽ các node mapping, assessment mức A, R-EXACT-LABEL,
hai premise nhãn chuẩn hóa và evidence; ghi tên cạnh theo RDF đã xuất.
Ảnh cần cho thấy đường đi từ assessment đến dữ kiện hỗ trợ.}
{fig:proof}
